\section{Supplementary proof details and scope audit}
\label{app:proofs}

\subsection{Fourier uniqueness in the classification theorem}

If a finite signed measure $\nu$ has all Fourier coefficients zero, then it
annihilates every trigonometric polynomial.  For $f\in C(\TT)$, choose
trigonometric polynomials $P_n$ with $\|f-P_n\|_\infty\to0$.  Finite total
variation gives
\[
 \left|\int f\,d\nu\right|
 \le \left|\int P_n\,d\nu\right|
     +\|f-P_n\|_\infty\,\|\nu\|_{\rm TV}\longrightarrow0.
\]
Thus $\nu$ annihilates $C(\TT)$ and is zero.  No moment-determinacy theorem
beyond compact Fourier uniqueness is needed.

\subsection{Faithfulness and normalization}

For $c>0$, if $a\oplus x\ge0$ and $\Phi_c(a\oplus x)=0$, then $a=0$ and
$\tau(x)=0$.  Faithfulness of the canonical group trace gives $x=0$.
Nevertheless $\Phi_c(1)=1+c$, so the word ``faithful'' does not imply
``state.''  This distinction is the price of the diffuse escape.

\subsection{Jensen boundary cases}

For $|q|<1$, the circle average in \cref{prop:fk} is zero.  For $|q|>1$,
factor $1-qz=-qz(1-q^{-1}z^{-1})$; the second factor has zero mean logarithm,
leaving $\log|q|$.  At $|q|=1$, the integrable logarithmic singularity follows
by radial limiting.  Hence the formula holds on all of $\CC$ with determinant
zero at $q=1$ because of the scalar point mass.

\subsection{What is not proved}

The classification does not rule out complex signed measures with the same
positive moments, non-scalar operator-valued determinants, or character
families before Haar averaging.  None of those variants may inherit the
Route-A tuple without a new source lock, determinant convention, and matched
arithmetic controls.  Cross-family geometric or operator constructions are
recorded only in \texttt{ROUND2\_CLUES.md}.
